% XeLaTeX; A2 design scaled to A4 for the preview/export.
\documentclass{article}
\usepackage[paperwidth=420mm,paperheight=594mm,margin=18mm]{geometry}
\usepackage{pgfpages}
\pgfpagesuselayout{resize to}[a4paper]
\usepackage{fontspec}
\setmainfont{BaiJamjuree}[Path=font/Bai_Jamjuree/,Extension=.ttf,UprightFont=*-Regular,BoldFont=*-Bold,ItalicFont=*-Italic,BoldItalicFont=*-BoldItalic]
\XeTeXlinebreaklocale "th"
\XeTeXlinebreakskip=0pt plus 1pt
\XeTeXgenerateactualtext=1
\usepackage{xcolor,graphicx,tikz,array,tabularx}
\usetikzlibrary{arrows.meta,positioning}
\definecolor{ink}{HTML}{153A3C}
\definecolor{teal}{HTML}{087F75}
\definecolor{mint}{HTML}{E9F5EF}
\definecolor{blue}{HTML}{3159A7}
\definecolor{pale}{HTML}{EEF3FC}
\definecolor{muted}{HTML}{506865}
\definecolor{paper}{HTML}{FAFCF8}
\pagestyle{empty}
\setlength{\parindent}{0pt}
\setlength{\parskip}{3mm}
\long\def\placeposterblock#1,#2\stop#3#4{%
  \AddToHook{shipout/foreground}{\setlength{\unitlength}{1mm}\put(#1,-#2){\raisebox{-\height}{
        \begin{minipage}[t]{#3mm}#4
\end{minipage}}}}}
\NewDocumentEnvironment{textblock}{m r() +b}{\placeposterblock#2\stop{#1}{#3}}{}
\newcommand{\heading}[2]{{\color{teal}\bfseries\fontsize{27}{35}\selectfont #1\quad #2\par}\vspace{2mm}}
\newcommand{\smalllabel}[1]{{\fontsize{16}{23}\selectfont\color{muted}#1\par}}
\newcommand{\role}[2]{\textbf{#1} & #2\\[4mm]}
\begin{document}
\color{ink}\fontsize{21}{31}\selectfont
\AddToHook{shipout/background}{
  \begin{tikzpicture}[remember picture,overlay]
    \fill[paper] (current page.south west) rectangle (current page.north east);
    \fill[teal] ([xshift=0mm,yshift=-9mm]current page.north west) rectangle ([xshift=420mm,yshift=-14mm]current page.north west);
    \fill[mint,rounded corners=5mm] ([xshift=20mm,yshift=-116mm]current page.north west) rectangle ([xshift=400mm,yshift=-159mm]current page.north west);
    \fill[white,rounded corners=5mm] ([xshift=20mm,yshift=-239mm]current page.north west) rectangle ([xshift=204mm,yshift=-535mm]current page.north west);
    \fill[pale,rounded corners=5mm] ([xshift=216mm,yshift=-239mm]current page.north west) rectangle ([xshift=400mm,yshift=-535mm]current page.north west);
    \draw[teal!25,line width=0.6pt] ([xshift=20mm,yshift=-545mm]current page.north west)--([xshift=400mm,yshift=-545mm]current page.north west);
\end{tikzpicture}}

\begin{textblock}{39}(23,25)
  \includegraphics[width=33mm,height=41mm,keepaspectratio]{images/logo.png}
\end{textblock}
\begin{textblock}{327}(70,25)
  \smalllabel{STUDENT LEAVE DATA WAREHOUSE}
  {\bfseries\fontsize{43}{54}\selectfont ข้อมูลลา เล่าอะไรได้บ้าง?\par}
  {\fontsize{24}{34}\selectfont คลังข้อมูลการลานิสิต คณะวิศวกรรมศาสตร์ กำแพงแสน\par}
\end{textblock}
\begin{textblock}{376}(24,84)
  {\fontsize{18}{27}\selectfont
    นายจิรภัทร แซ่ลี้ 6620501311\quad นายกานตภัท อาจหาญ 6620503349\\
  นายธนวัฒน์ แซ่จึง 6620503403\quad นายศุภกร ลีลาธนางกุล 6620503578\par}
\end{textblock}
\begin{textblock}{359}(31,123)
  {\bfseries\fontsize{25}{34}\selectfont จากใบลาหลายรายการ → สู่ภาพรวมที่เข้าใจง่าย\par}
  รวมข้อมูลให้ดูได้ทั้งรายคน รายวิชา และรายเทอม\quad เลือกมุมมองตามบทบาท\\
  โดยแยกหน่วยนับให้ชัด เพื่อไม่ให้ยอดซ้ำเมื่อใบลาเดียวเกี่ยวข้องกับหลายวิชา
\end{textblock}

\begin{textblock}{376}(24,171)
  \heading{01}{ข้อมูลเดินทางมายังไง}
  \begin{tikzpicture}[>=Stealth,box/.style={rounded corners=3mm,draw=teal!40,fill=white,text width=77mm,minimum height=24mm,align=center,font=\fontsize{20}{28}\selectfont},node distance=15mm]
    \node[box] (a) {ข้อมูลใบลา\\และข้อมูลอ้างอิง};
    \node[box,right=of a] (b) {ETL\\ตรวจและจัดรูปข้อมูล};
    \node[box,right=of b] (c) {Warehouse\\เก็บเป็นโครงสร้างกลาง};
    \node[box,right=of c] (d) {รายงานตามบทบาท\\กราฟ ตาราง รายละเอียด};
    \draw[->,thick,teal](a)--(b);\draw[->,thick,teal](b)--(c);\draw[->,thick,teal,dashed](c)--(d);
  \end{tikzpicture}
  \smalllabel{เส้นประ: การเชื่อมรายงานกับข้อมูลจริงยังต้องตรวจครบเส้นทาง หน้าเว็บที่ตรวจปัจจุบันใช้ข้อมูลจำลอง}
\end{textblock}

\begin{textblock}{166}(29,248)
  \heading{02}{เก็บแยกกัน เพื่อรวมให้ถูก}
  \textbf{Fact ใบลา}\quad 1 แถว = 1 ใบลา\\
  ใช้สรุปจำนวนใบลา วันลา และผู้ที่ลา

  \textbf{Fact ใบลา–กลุ่มเรียน}\\
  1 แถว = 1 ใบลา × 1 กลุ่มเรียน\\
  ใช้ดูผลกระทบที่เกี่ยวกับแต่ละวิชา

  \vspace{1mm}
  \begin{tikzpicture}[>=Stealth,box/.style={draw=teal!40,fill=mint,rounded corners=2mm,align=center,text width=68mm,minimum height=18mm,font=\fontsize{19}{27}\selectfont}]
    \node[box] (dim) at (0,0) {มิติข้อมูล\\นิสิต · เทอม · ประเภทลา};
    \node[box] (req) at (0,-2.7) {Fact ใบลา\\จำนวนใบ · วันลา};
    \node[box] (org) at (8.1,0) {มิติข้อมูล\\หน่วยงาน · ที่ปรึกษา};
    \node[box] (rs) at (8.1,-2.7) {Fact ใบลา–กลุ่มเรียน\\ครั้งเรียน · ชั่วโมงเรียน};
    \draw[->,thick,teal](dim)--(req);\draw[->,thick,teal](org)--(rs);
    \draw[->,thick,teal](req)--node[above,font=\fontsize{16}{23}\selectfont]{1 : หลาย}(rs);
  \end{tikzpicture}
  \smalllabel{รายวิชา กลุ่มเรียน และตารางเชื่อมผู้สอนช่วยให้รองรับหลายวิชาและผู้สอนร่วม}

  \vspace{3mm}
  \heading{03}{ใบเดียว หลายวิชา ก็ไม่ซ้ำ}
  \textbf{สมมติลา 2 วัน เกี่ยวข้องกับ 3 วิชา}\\
  ภาพรวม = \textbf{1 ใบ / 2 วัน}\\
  รายละเอียดวิชา = \textbf{3 ความสัมพันธ์}\\
  จึงไม่เอา 2 × 3 ไปนับเป็น 6 วันลา

  \smalllabel{วันลาตามปฏิทิน ≠ ชั่วโมงเรียนที่ได้รับผลกระทบ\\ชั่วโมงเรียนต้องมีข้อมูลตารางเรียนครบก่อน}
\end{textblock}

\begin{textblock}{166}(225,248)
  \heading{04}{แต่ละคนอยากรู้อะไร}
  \renewcommand{\arraystretch}{1.1}
  \begin{tabularx}{\linewidth}{@{}p{48mm}X@{}}
    \role{นิสิต}{เทอมนี้เราลาเท่าไร\newline เกี่ยวข้องกับวิชาไหนบ้าง}
    \role{ที่ปรึกษา}{ในกลุ่มที่ดูแล ใครหรือ\newline ชั้นปีไหนมีการลาสูง}
    \role{ผู้สอน}{วิชาและกลุ่มเรียนที่สอน\newline มีการลามากช่วงไหน}
    \role{หัวหน้าภาค}{นิสิตในภาคกับวิชาของภาค\newline มีภาพรวมต่างกันยังไง}
    \role{ผู้บริหาร}{แต่ละภาคและหลักสูตร\newline มีรูปแบบการลาต่างกันไหม}
    \role{เจ้าหน้าที่}{กลุ่มนิสิตหรือที่ปรึกษาไหน\newline ควรเปิดดูรายละเอียดต่อ}
  \end{tabularx}
  \smalllabel{แต่ละบทบาทเห็นข้อมูลตามขอบเขตที่ได้รับสิทธิ์\\ไม่ใช้รหัสใน URL เป็นสิทธิ์เปิดข้อมูล}

  \vspace{4mm}
  \heading{05}{อ่านกราฟให้ได้คำตอบ}
  เลือก \textbf{ปี / เทอม / ประเภทลา} ให้ตรงกัน\\
  ค่อยเทียบจำนวนใบ วันลา หรือรายวิชา

  \textbf{แนวโน้ม} → เปลี่ยนไปยังไงในแต่ละเทอม\\
  \textbf{Top 10} → เริ่มดูรายละเอียดจากกลุ่มไหน\\
  \textbf{สัดส่วน} → แต่ละกลุ่มมีส่วนเท่าไร

  \smalllabel{ยอดสูงยังไม่บอกสาเหตุหรือความเสี่ยงทันที\\ต้องดูจำนวนนิสิตและบริบทของแต่ละกลุ่มด้วย}
\end{textblock}

\begin{textblock}{376}(24,552)
  {\bfseries\fontsize{19}{28}\selectfont ไปต่อยังไง?\quad เชื่อมข้อมูลจริง → ตรวจยอดและสิทธิ์ → ทดลองใช้กับผู้ใช้\par}
  \smalllabel{สถานะ: สรุปจากโครงสร้างและโค้ด ไม่ใช่ผลประเมินการใช้งานจริง ไม่มีการอ้างตัวเลขลดเวลาหรือคะแนนความพึงพอใจ}
  {\fontsize{16}{23}\selectfont\color{muted}ภาควิชาวิศวกรรมคอมพิวเตอร์ · คณะวิศวกรรมศาสตร์ กำแพงแสน · มหาวิทยาลัยเกษตรศาสตร์\par}
\end{textblock}
\null
\end{document}
